% =====================================================================
%  2 Fundamentação Teórica
% =====================================================================
\section{Fundamentação Teórica}
\label{sec:fundamentacao}

Esta seção apresenta os conceitos que sustentam a análise: a Engenharia
Semiótica e o Método de Inspeção Semiótica, os papéis que as pessoas
assumem em casas inteligentes compartilhadas, as rotinas como forma de
desenvolvimento pelo usuário final e a noção de integridade contextual.

\subsection{Engenharia Semiótica e metacomunicação}
\label{sec:engsem}

A Engenharia Semiótica é uma teoria explicativa da IHC que trata a
interface de um sistema interativo como uma comunicação enviada por seu
designer aos usuários~\cite{desouza2005semiotic}. Essa comunicação é
unidirecional: o designer não está presente durante o uso e fala por meio
da interface, que os usuários interpretam enquanto interagem com
ela~\cite{fortunato2025ihc}. Por se tratar de uma comunicação sobre como
se comunicar com o sistema, a teoria a denomina metacomunicação.

O conteúdo da metamensagem pode ser parafraseado por um modelo em que o
designer diz ao usuário qual é o seu entendimento de quem ele é, do que
ele quer ou precisa fazer, de que maneiras prefere fazê-lo e por quê; e
que, por isso, aquele é o sistema projetado para ele, que pode ou deve
ser usado de determinada forma para atingir um conjunto de objetivos
compatíveis com essa visão~\cite[p.~25]{desouza2005semiotic}. Esse modelo
se dirige a um interlocutor, o ``você'' a quem o designer fala. Em
sistemas de uso individual, esse interlocutor coincide com a pessoa que
usa o sistema. Em um dispositivo instalado em um espaço coletivo, cabe
perguntar quem é esse ``você'', questão retomada na
Seção~\ref{sec:multiusuario}.

A metamensagem é composta por signos, classificados em três
tipos~\cite{desouza2006sim,desouza2009methods}. Os \textbf{signos
estáticos} expressam o estado do sistema e podem ser interpretados sem
relações causais ou temporais, considerando apenas os elementos presentes
na interface em um dado momento, como rótulos, estados de botões e opções
selecionadas. Os \textbf{signos dinâmicos} expressam o comportamento do
sistema e só podem ser percebidos por meio da interação, como o efeito de
acionar um botão. Os \textbf{signos metalinguísticos} referem-se a outros
signos da interface e são usados pelo designer para explicar o que eles
significam e como usá-los, como mensagens de ajuda, dicas e mensagens de
erro~\cite{fortunato2025ihc}.

A qualidade de um sistema interativo, nessa perspectiva, é a sua
comunicabilidade: a capacidade de comunicar de forma organizada e
consistente a lógica, a intenção e os princípios de
design~\cite{desouza2005semiotic,desouza2009methods}. A teoria propõe dois
métodos principais para avaliá-la. O Método de Avaliação de
Comunicabilidade (MAC) observa usuários reais e avalia como a
metamensagem é recebida. O Método de Inspeção Semiótica (MIS) é conduzido
por especialistas e avalia como a metamensagem é
emitida~\cite{desouza2009methods,fortunato2025ihc}. Por não envolver
participantes, o MIS é o método adotado neste trabalho.

\subsection{O Método de Inspeção Semiótica e suas variantes}
\label{sec:mis}

No MIS, o avaliador analisa o que cada tipo de signo comunica,
reconstrói a metamensagem do designer e identifica potenciais rupturas de
comunicação~\cite{desouza2006sim}. O método é organizado nas seguintes
etapas~\cite{desouza2006sim,fortunato2025ihc}:

\begin{itemize}
  \item \textbf{Preparação:} inspeção informal, definição do foco da
        análise, dos usuários pretendidos, dos objetivos apoiados pelo
        sistema e dos cenários de inspeção;
  \item \textbf{Etapas 1, 2 e 3:} inspeção dos signos metalinguísticos,
        estáticos e dinâmicos, respectivamente, com a reconstrução de uma
        versão da metamensagem e a identificação de potenciais rupturas a
        partir de cada tipo de signo;
  \item \textbf{Etapa 4:} contraste das três versões da metamensagem, em
        busca de significados contraditórios ou ambíguos;
  \item \textbf{Etapa 5:} apreciação final da comunicabilidade do sistema,
        registrada em um relatório.
\end{itemize}

Durante a inspeção, o avaliador atua como representante dos usuários e
deve ser especialista em IHC e conhecer a Engenharia
Semiótica~\cite{desouza2006sim,fortunato2025ihc}.

Uma revisão sistemática recente identificou cinco variantes do método,
que, somadas à versão original, formam seis versões~\cite{fortunato2025ihc}.
A \textbf{MIS Científica}~\cite{desouza2010scientific} acrescenta uma
questão de pesquisa na preparação e uma etapa de triangulação dos
resultados com evidências de outras fontes, como materiais publicados na
web sobre o sistema, entre eles opiniões de usuários e perguntas
frequentes~\cite{fortunato2025ihc}. A MIS com signos dêiticos inclui uma
nova classe de signos, que se refere a elementos externos à interface: a
pessoa (identidade do usuário), o espaço e o tempo da
interação~\cite{fortunato2025ihc}. As demais variantes inspecionam mais de
uma versão da metamensagem: a CP-SIM, entre plataformas de um mesmo
sistema; a MIS mediada por leitor de tela
(SIM-SR)~\cite{carvalho2023simsr}, entre a interface original e a
interface percebida por meio do leitor de tela; e a SIM-HInt, entre os
componentes de sistemas de integração humano-computador~\cite{fortunato2025ihc}.

As variantes ampliam o método na direção de novos tipos de signos e de
múltiplas versões ou componentes de um sistema. Nenhuma das cinco, tal
como descritas por \citet{fortunato2025ihc}, tem como foco a existência de
mais de um destinatário da metamensagem, que é a situação examinada neste
trabalho.

\subsection{Casas inteligentes com múltiplos usuários}
\label{sec:multiusuario}

Diferentemente de computadores e smartphones, dispositivos de casa
inteligente instalados em espaços compartilhados tornam-se dispositivos
usados por várias pessoas e que afetam todas elas~\cite{geeng2019control}.
Os estudos sobre esses ambientes descrevem papéis distintos entre os
moradores. \citet{geeng2019control} usam o termo \textit{smart home
driver} para quem toma a iniciativa de adotar a tecnologia, sobretudo
instalando os dispositivos, e \textit{passive user} para quem participa
menos dessas decisões. \citet{koshy2021passenger} distinguem os usuários
\textit{pilotos}, que introduzem novas funcionalidades na casa, instalando
dispositivos ou criando rotinas, dos usuários \textit{passageiros}, que
usam o que já existe. Em seu levantamento, pelo menos 52\% dos
respondentes que moravam com outras pessoas viviam em casas com ao menos
um piloto e um passageiro~\cite{koshy2021passenger}.

Esses estudos convergem em alguns achados. Há uma concentração de
conhecimento, acesso e controle na pessoa que seleciona e instala os
dispositivos, que frequentemente também é a única com acesso à conta à
qual eles estão vinculados~\cite{geeng2019control}. Muitos
\textit{drivers} não consultam os demais moradores antes de instalar
dispositivos, e estes passam a depender deles quando algo deixa de
funcionar~\cite{geeng2019control}. O conhecimento sobre os dispositivos
flui dos pilotos para os passageiros, e a casa tende a refletir os valores
dos pilotos, o que dificulta que os passageiros incorporem os dispositivos
ao seu cotidiano~\cite{koshy2021passenger}. Em um dos domicílios
estudados por \citet{koshy2021passenger}, o piloto praticamente deixou de
criar rotinas porque elas incomodavam a esposa: as rotinas matinais, com
horário fixo e pensadas para a agenda dele, não combinavam com o horário
de trabalho dela. \citet{zeng2017concerns} observaram
um descompasso semelhante: os moradores que não administram a casa
tendem a estar menos informados sobre os riscos de segurança e
privacidade. \citet{zeng2019multiuser} mostram que esse desequilíbrio de
poder aparece também em relações não abusivas e propõem como princípios
de design a flexibilidade no controle de acesso, a agência de todos os
usuários, o respeito entre eles e a transparência do comportamento da
casa.

Além dos moradores, há pessoas que convivem com os dispositivos sem
participar de sua configuração. \citet{bernd2020bystanders} estudaram
babás que trabalham em casas com dispositivos de vigilância e as
analisaram como pessoas afetadas (\textit{bystanders}), em uma relação
marcada pela diferença de poder entre empregador e empregada.
\citet{geeng2019control} relatam uma tensão semelhante envolvendo o acesso
de uma diarista a uma fechadura inteligente.

Os assistentes de voz acentuam o caráter coletivo desses dispositivos. O
uso da Amazon Echo se entrelaça com as conversas e as atividades
simultâneas dos moradores, como as refeições em
família~\cite{porcheron2018voice}, e a interface de voz, acessível a todos,
amplia o acesso da família à tecnologia~\cite{beneteau2020parenting}. A
configuração do dispositivo e de suas rotinas, no entanto, continua
concentrada no aplicativo vinculado à conta de uma
pessoa~\cite{barricelli2024routines}.

\subsection{Rotinas como desenvolvimento pelo usuário final}
\label{sec:rotinas}

Rotinas são sequências de ações definidas pelo usuário e disparadas por
eventos ou condições específicas~\cite{barricelli2022multimodal}. Na
Amazon Alexa e no Google Nest, elas são criadas por meio do aplicativo
complementar no smartphone~\cite{barricelli2024routines}. Essa forma de
configuração corresponde à programação gatilho-ação (``se \ldots{},
então \ldots{}''), na qual \citet{ur2014tap} mostraram que é possível
expressar a maior parte dos comportamentos desejados pelos participantes
de seu estudo e que usuários sem experiência aprendem rapidamente a criar
regras com múltiplos gatilhos ou ações.

Criar a rotina, porém, é só uma parte do uso. \citet{koshy2021passenger},
retomando o estudo longitudinal de \citet{woo2015diy}, relatam que os
usuários dedicam tempo a refinar rotinas depois de criadas, abandonam
muitas delas e que os autores sugeriram oferecer um canal para que os
moradores avaliem as rotinas configuradas por outros membros da casa.
Antes da popularização dos assistentes de voz, \citet{brush2011wild} já
apontavam a pouca flexibilidade e a dificuldade de gerenciar a automação
residencial como barreiras à sua adoção.

\subsection{Integridade contextual}
\label{sec:integridade}

Para analisar o uso de dados nas rotinas, adotamos a noção de integridade
contextual. \citet{nissenbaum2004privacy} propõe que a proteção da
privacidade depende das normas de cada contexto: a coleta e a circulação
de informações devem ser adequadas ao contexto e obedecer às normas de
distribuição que o regem. A autora distingue as normas de adequação, que
definem quais informações sobre uma pessoa é apropriado revelar em um
contexto, das normas de fluxo ou distribuição, que definem como essas
informações podem circular, e afirma que não há esfera da vida em que
``tudo é permitido''. \citet{bernd2020bystanders} usam esse referencial
para analisar como as relações de poder entre moradores e pessoas afetadas
moldam as normas de privacidade em casas inteligentes. No caso das
rotinas, um dado associado a uma pessoa, como o horário de seu alarme,
pode disparar ações percebidas por todos que estão no ambiente, o que
caracteriza uma mudança no fluxo dessa informação.
